\documentclass[11pt,a4paper]{article}

\usepackage[utf8]{inputenc}
\usepackage[T1]{fontenc}
\usepackage[provide=*,english]{babel}
\usepackage{lmodern}
\usepackage{microtype}
\usepackage{amsmath,amssymb,mathtools,amsthm}
\usepackage{geometry}
\usepackage{xcolor}
\usepackage{enumitem}
\usepackage{booktabs}
\usepackage{fancyhdr}
\usepackage[hidelinks]{hyperref}

\geometry{left=2.3cm,right=2.3cm,top=2.2cm,bottom=2.25cm,headheight=14pt}
\setlength{\parindent}{0pt}
\setlength{\parskip}{0.38em}
\setlist[itemize]{topsep=0.25em,itemsep=0.18em,leftmargin=1.5em}
\setlist[enumerate]{topsep=0.3em,itemsep=0.25em,leftmargin=1.8em}
\allowdisplaybreaks

\definecolor{blue}{RGB}{34,73,110}
\definecolor{gray}{RGB}{90,96,102}

\pagestyle{fancy}
\fancyhf{}
\fancyhead[L]{\small\color{gray}GC2D: symplectic integration}
\fancyhead[R]{\small\color{gray}Generalized-energy near-conservation}
\fancyfoot[C]{\thepage}
\renewcommand{\headrulewidth}{0.3pt}

\newtheorem{theorem}{Theorem}[section]
\newtheorem{lemma}[theorem]{Lemma}
\newtheorem{proposition}[theorem]{Proposition}
\newtheorem{corollary}[theorem]{Corollary}
\theoremstyle{definition}
\newtheorem{assumption}[theorem]{Assumption}
\theoremstyle{remark}
\newtheorem{remark}[theorem]{Remark}

\newcommand{\R}{\mathbb{R}}
\newcommand{\J}{\mathbb{J}}
\newcommand{\Om}{\Omega}
\newcommand{\K}{\mathcal{K}}
\newcommand{\Kh}{\widetilde{\mathcal{K}}_h}
\newcommand{\PhiNum}{\Phi_h}
\newcommand{\flow}[1]{\varphi^{#1}}
\newcommand{\norm}[1]{\left\lVert #1\right\rVert}
\newcommand{\abs}[1]{\left\lvert #1\right\rvert}

\begin{document}

\begin{center}
  {\LARGE\bfseries\color{blue}Absence of Secular Generalized-Energy Drift}\par
  \vspace{0.35em}
  {\large A backward-error argument for Gauss--Legendre and fully extended ABBA methods}\par
  \vspace{0.45em}
  {\small Precise theorem, proof, scope, and implementation limitations for GC2D}
\end{center}

\vspace{0.5em}

\section{Purpose and precise claim}

This note proves a conditional near-conservation result for fixed-step
symplectic discretizations of the time-dependent guiding-centre problem.  The
central conclusion is not exact conservation and is not a statement for
infinite time.  Under the analytic backward-error hypotheses stated below, a
symplectic method of order $p$ satisfies
\begin{equation}
 \abs{\K(Z_n)-\K(Z_0)}\leq C h^p
 \qquad\text{for}\qquad
 t_n=nh\leq \exp(c/h),
 \label{eq:headline-bound}
\end{equation}
after a harmless reduction of the positive constant $c$.  The constants are
independent of $h$ and $n$, but they depend on the Hamiltonian, the complex
analyticity domain, and the compact set containing the numerical orbit.

Equation \eqref{eq:headline-bound} excludes a secular contribution of the
form $t_nh^p$ throughout the stated interval.  It does not imply that the
energy error is exactly periodic, that its average is zero, or that the same
bound holds as $n\to\infty$ with $h$ fixed.

The result applies to the autonomous extended Hamiltonian
\begin{equation}
 \K(x,y,t,k)=H(t,x,y)+k,
 \label{eq:extended-hamiltonian-introduction}
\end{equation}
not to the explicitly time-dependent physical value $H(t,x,y)$ alone.

\section{Physical and extended Hamiltonian systems}

For one guiding-centre particle, let
\[
 z=(x,y)^T,
 \qquad
 J_0=
 \begin{pmatrix}
  0&-1\\
  1&0
 \end{pmatrix}.
\]
The physical system is
\begin{equation}
 \dot z=J_0\nabla_zH(t,z)
 =
 \begin{pmatrix}
  -H_y(t,x,y)\\
  H_x(t,x,y)
 \end{pmatrix}.
 \label{eq:physical-hamiltonian-system}
\end{equation}
Along an exact physical trajectory,
\begin{align}
 \frac{d}{dt}H(t,z(t))
 &=H_t(t,z)+\nabla_zH(t,z)^T\dot z
 \notag\\
 &=H_t(t,z)+\nabla_zH(t,z)^TJ_0\nabla_zH(t,z)
 =H_t(t,z).
 \label{eq:physical-energy-change}
\end{align}
The last equality uses the skew-symmetry of $J_0$.  Consequently, $H$ is not
an invariant when $H_t\neq0$.

Introduce the extended state and constant Poisson matrix
\begin{equation}
 Z=(x,y,t,k)^T,
 \qquad
 \J=
 \begin{pmatrix}
  0&-1&0&0\\
  1&0&0&0\\
  0&0&0&1\\
  0&0&-1&0
 \end{pmatrix}.
 \label{eq:extended-poisson-matrix}
\end{equation}
With $\K$ from \eqref{eq:extended-hamiltonian-introduction}, the autonomous
extended equations are
\begin{equation}
 \dot Z=\J\nabla_Z\K(Z)
 =
 \begin{pmatrix}
  -H_y\\ H_x\\ 1\\ -H_t
 \end{pmatrix}.
 \label{eq:extended-system}
\end{equation}
They reproduce the physical system, advance time at unit speed, and evolve
the conjugate momentum according to $\dot k=-H_t$.  Moreover,
\begin{equation}
 \frac{d}{ds}\K(Z(s))
 =\nabla_Z\K(Z)^T\J\nabla_Z\K(Z)=0.
 \label{eq:exact-generalized-energy-conservation}
\end{equation}
Thus $\K=H+k$ is the invariant whose numerical drift is meaningful.

Since $\J^{-1}=-\J$, the extended symplectic-form matrix is
\begin{equation}
 \Om=\J^{-1}=-\J.
 \label{eq:extended-symplectic-form}
\end{equation}
A differentiable numerical map $\PhiNum$ is symplectic when
\begin{equation}
 D\PhiNum(Z)^T\Om D\PhiNum(Z)=\Om.
 \label{eq:symplectic-map-definition}
\end{equation}

\section{Hypotheses for exponentially long near-conservation}

The exponential time scale requires stronger assumptions than finite-interval
convergence.

\begin{assumption}[Analytic Hamiltonian]
\label{ass:analyticity}
There is a compact simply connected real domain $D\subset\R^4$ and a complex
neighbourhood $D_\rho$ of uniform width $\rho>0$ on which $\K$ is analytic
and bounded together with the derivatives required below.
\end{assumption}

\begin{assumption}[Ideal numerical map]
\label{ass:ideal-map}
For fixed $h$, the one-step map $\PhiNum$ is analytic on $D_\rho$, symplectic,
and of order $p$:
\begin{equation}
 \PhiNum(Z)=\flow{\K}_h(Z)+O(h^{p+1}),
 \label{eq:method-order}
\end{equation}
uniformly on the relevant domain.  Every implicit stage or projection equation
is solved at its exact local root and arithmetic is exact.
\end{assumption}

\begin{assumption}[Confined orbit and small fixed step]
\label{ass:confined-orbit}
The fixed step $h$ is sufficiently small relative to the analytic scales of
the problem, and the numerical orbit $Z_j=\PhiNum^j(Z_0)$ remains in a compact
subset of $D$ for all indices considered.
\end{assumption}

The qualification ``sufficiently small'' is problem dependent.  For a
Hamiltonian containing a fast frequency $\omega$, the constants deteriorate
as $\omega$ increases; resolving that scale, rather than merely requiring a
small dimensional value of $h$, is essential.

\section{Modified Hamiltonian}

Backward-error analysis seeks a differential equation whose exact time-$h$
flow reproduces the numerical map to a much higher accuracy than the original
order estimate.  Formally, the modified vector field has an expansion
\begin{equation}
 F_h=F+h^pF_p+h^{p+1}F_{p+1}+\cdots,
 \qquad
 F=\J\nabla\K.
 \label{eq:formal-modified-vector-field}
\end{equation}

\begin{proposition}[Hamiltonian character of the modified equation]
\label{prop:modified-hamiltonian}
Under Assumption~\ref{ass:ideal-map}, every coefficient of the formal modified
vector field is locally Hamiltonian.  On the simply connected domain $D$, one
may therefore write
\begin{equation}
 F_h=\J\nabla\Kh,
 \qquad
 \Kh
 =\K+h^p\K_p+h^{p+1}\K_{p+1}+\cdots.
 \label{eq:formal-modified-hamiltonian}
\end{equation}
If the method is also symmetric,
\begin{equation}
 \Phi_{-h}=\PhiNum^{-1},
 \label{eq:symmetry-condition}
\end{equation}
then the modified Hamiltonian may be chosen with the parity expansion
\begin{equation}
 \Kh=\K+h^p\K_p+h^{p+2}\K_{p+2}+h^{p+4}\K_{p+4}+\cdots.
 \label{eq:symmetric-modified-hamiltonian}
\end{equation}
\end{proposition}

\begin{proof}
The formal logarithm of a near-identity one-step map defines the modified
vector field.  Symplecticity of $\PhiNum$ implies, order by order, that its
formal logarithm is a symplectic vector field.  On a simply connected canonical
domain, a symplectic vector field is Hamiltonian, which gives
\eqref{eq:formal-modified-hamiltonian}.  The order-$p$ consistency in
\eqref{eq:method-order} removes all corrections below $h^p$.

For a symmetric method, replacing $h$ by $-h$ produces the inverse map.
Taking the formal logarithm shows that the modified vector field is an even
function of $h$.  Hence the terms after $h^p$ advance in increments of two,
as stated in \eqref{eq:symmetric-modified-hamiltonian}.  Symmetry supplies this
parity, but symplecticity is the property that makes the modified vector field
Hamiltonian.
\end{proof}

The formal series generally diverges and cannot itself be declared an exactly
conserved numerical energy.  Analyticity permits optimal truncation before
the coefficients begin to grow.

\begin{lemma}[Exponentially accurate truncated modified Hamiltonian]
\label{lem:exponential-modified-hamiltonian}
Under Assumptions~\ref{ass:analyticity}--\ref{ass:confined-orbit}, there are
positive constants $C_0$, $C_1$, $c$, and $h_0$ such that, for
$0<h<h_0$, the formal series can be truncated at an index $N(h)=O(h^{-1})$
to produce $\Kh^{[N]}$ satisfying
\begin{align}
 \sup_{Z\in D}\abs{\Kh^{[N]}(Z)-\K(Z)}
 &\leq C_0h^p,
 \label{eq:modified-original-distance}\\
 \abs{\Kh^{[N]}(\PhiNum(Z))-\Kh^{[N]}(Z)}
 &\leq C_1h\exp(-c/h)
 \label{eq:one-step-modified-energy-defect}
\end{align}
whenever the step and the connecting modified trajectory remain in the
domain used for the estimates.
\end{lemma}

\begin{proof}[Proof sketch]
Analyticity on the uniform complex neighbourhood gives Cauchy estimates for
successive derivatives.  Applied to the recursive construction of the
modified equation, these estimates bound its coefficients by factorially
growing quantities.  Truncating near the smallest term, at
$N(h)\asymp c/h$, converts the algebraic remainder into an exponentially small
one.  Because Proposition~\ref{prop:modified-hamiltonian} makes the truncated
modified equation Hamiltonian, its exact flow conserves $\Kh^{[N]}$.  The
exponentially small difference between that flow and $\PhiNum$ then gives
\eqref{eq:one-step-modified-energy-defect}; consistency gives
\eqref{eq:modified-original-distance}.  This optimal-truncation result is the
standard analytic input of Hamiltonian backward-error analysis.
\end{proof}

\section{No-secular-drift theorem}

\begin{theorem}[Generalized-energy near-conservation]
\label{thm:no-secular-drift}
Under Assumptions~\ref{ass:analyticity}--\ref{ass:confined-orbit}, there are
positive constants $C$, $c$, and $h_0$ such that the numerical orbit satisfies
\begin{equation}
 \abs{\K(Z_n)-\K(Z_0)}
 \leq Ch^p+C t_n\exp(-c/h),
 \qquad t_n=nh,
 \label{eq:full-energy-bound}
\end{equation}
for every index for which the domain assumptions remain valid.  Consequently,
after possibly reducing $c$, there is a constant $C'$ such that
\begin{equation}
 \boxed{
 \abs{\K(Z_n)-\K(Z_0)}\leq C'h^p
 \qquad\text{when}\qquad
 0\leq t_n\leq\exp\!\left(\frac{c}{2h}\right).}
 \label{eq:exponentially-long-bound}
\end{equation}
\end{theorem}

\begin{proof}
Let $\Kh^{[N]}$ be supplied by
Lemma~\ref{lem:exponential-modified-hamiltonian}.  Telescoping its one-step
defect along the numerical orbit gives
\begin{align}
 \abs{\Kh^{[N]}(Z_n)-\Kh^{[N]}(Z_0)}
 &\leq
 \sum_{j=0}^{n-1}
 \abs{\Kh^{[N]}(Z_{j+1})-\Kh^{[N]}(Z_j)}
 \notag\\
 &\leq nC_1h\exp(-c/h)
 =C_1t_n\exp(-c/h).
 \label{eq:telescoped-modified-energy}
\end{align}
Insert and subtract the truncated modified Hamiltonian at both endpoints:
\begin{align}
 \abs{\K(Z_n)-\K(Z_0)}
 &\leq
 \abs{\K(Z_n)-\Kh^{[N]}(Z_n)}
 \notag\\
 &\quad+
 \abs{\Kh^{[N]}(Z_n)-\Kh^{[N]}(Z_0)}
 +\abs{\Kh^{[N]}(Z_0)-\K(Z_0)}.
 \label{eq:energy-triangle-inequality}
\end{align}
Equations \eqref{eq:modified-original-distance} and
\eqref{eq:telescoped-modified-energy} yield
\[
 \abs{\K(Z_n)-\K(Z_0)}
 \leq2C_0h^p+C_1t_n\exp(-c/h),
\]
which is \eqref{eq:full-energy-bound} after renaming constants.

If $t_n\leq\exp(c/(2h))$, the accumulated exponential term is bounded by
$C_1\exp(-c/(2h))$.  For every fixed $p$, this is smaller than a constant
multiple of $h^p$ when $h$ is sufficiently small.  Absorbing it into the
first term proves \eqref{eq:exponentially-long-bound}.
\end{proof}

\begin{remark}[Meaning of ``no secular drift'']
Theorem~\ref{thm:no-secular-drift} rules out an energy error growing like
$t_nh^p$ on the exponentially long interval.  It proves confinement to an
$O(h^p)$ band, not literal oscillation about zero.  A plotted sign change or
small fitted slope is useful evidence, but it is not part of the theorem.
\end{remark}

\begin{remark}[Not an infinite-time theorem]
For fixed $h$, the upper time in \eqref{eq:exponentially-long-bound} is finite.
The theorem therefore does not exclude later effects caused by resonances,
loss of confinement, exponentially small remainders, or arithmetic errors.
\end{remark}

\section{Application to the GC2D integrators}

\subsection{Two-stage Gauss--Legendre}

Let $A=(a_{ij})$, $b=(b_i)$, and $c=(c_i)$ be the two-stage Gauss--Legendre
coefficients.  Applying this Runge--Kutta method to
\eqref{eq:extended-system} gives the stage times
\begin{equation}
 T_i=t_n+h\sum_j a_{ij}=t_n+c_ih,
 \label{eq:gauss-stage-times}
\end{equation}
because $A\mathbf{1}=c$.  The physical stages solve
\begin{equation}
 Z_i^{\mathrm{phys}}
 =z_n+h\sum_j a_{ij}
 J_0\nabla_zH(T_j,Z_j^{\mathrm{phys}}),
 \label{eq:gauss-physical-stages}
\end{equation}
and the accepted physical and momentum updates are
\begin{align}
 z_{n+1}
 &=z_n+h\sum_i b_iJ_0\nabla_zH(T_i,Z_i^{\mathrm{phys}}),
 \label{eq:gauss-physical-update}\\
 k_{n+1}
 &=k_n-h\sum_i b_iH_t(T_i,Z_i^{\mathrm{phys}}).
 \label{eq:gauss-momentum-update}
\end{align}
The internal $k$ stages need not be solved because neither the physical field
nor $-H_t$ depends on $k$.  Thus the triangular implementation of
\eqref{eq:gauss-momentum-update} is algebraically the Gauss method on the full
extended system.  This is the extended map represented by
\texttt{GaussLegendre4(track\_energy=True)} in the project implementation.

Gauss coefficients satisfy the symplectic Runge--Kutta identities
\begin{equation}
 b_i a_{ij}+b_j a_{ji}=b_ib_j,
 \qquad i,j=1,2.
 \label{eq:gauss-symplectic-condition}
\end{equation}
The method is symmetric and has order four.

\begin{corollary}[Ideal Gauss--Legendre energy band]
\label{cor:gauss-energy-band}
If Assumptions~\ref{ass:analyticity} and \ref{ass:confined-orbit} hold and the
Gauss stage equations are solved exactly, then
\begin{equation}
 \abs{\K(Z_n)-\K(Z_0)}=O(h^4)
 \label{eq:gauss-energy-band}
\end{equation}
through the exponentially long interval of
Theorem~\ref{thm:no-secular-drift}.
\end{corollary}

\subsection{Fully extended implicit ABBA}

The fully extended ABBA construction duplicates the complete canonical state
$Z=(x,y,t,k)$, applies a composition of Hamiltonian shears on the doubled
space, and returns to the diagonal through an implicit symmetric projection.
The exact projected map is symplectic under the following local hypotheses:
\begin{itemize}
 \item every doubled shear is evaluated exactly;
 \item the projection root is locally unique;
 \item the residual Jacobian with respect to the multiplier is nonsingular;
 \item the projection equation is solved exactly.
\end{itemize}
The physical ABBA2 projection proof and its local hypotheses are developed in
the model note \texttt{ABBA2\_implicit.tex}, located under
\path{docs/models/abba2-implicit/}.  The fully extended construction applies
the same symmetric-projection mechanism to the complete canonical state.
Palindromic compositions of the projected second-order map remain symplectic,
and the single-projection fourth-order variant applies the same mechanism
around the complete symplectic base composition.

\begin{corollary}[Ideal fully extended ABBA energy bands]
\label{cor:abba-energy-bands}
Under the local exact-projection hypotheses above and
Assumptions~\ref{ass:analyticity} and \ref{ass:confined-orbit}, the fully
extended implicit ABBA methods satisfy the bounds in Table~\ref{tab:energy-bands}.
\end{corollary}

\begin{table}[ht]
 \centering
 \caption{Ideal generalized-energy bands from symplectic backward-error analysis.}
 \label{tab:energy-bands}
 \begin{tabular}{@{}lcc@{}}
  \toprule
  Method & Order $p$ & Energy band \\
  \midrule
  Two-stage Gauss--Legendre & $4$ & $O(h^4)$ \\
  \texttt{ABBA2Implicit}, fully extended & $2$ & $O(h^2)$ \\
  \texttt{ABBA4Implicit}, fully extended & $4$ & $O(h^4)$ \\
  \texttt{ABBA4ImplicitSingleProjection}, fully extended & $4$ & $O(h^4)$ \\
  \texttt{ABBA6Implicit}, fully extended & $6$ & $O(h^6)$ \\
  \bottomrule
 \end{tabular}
\end{table}

The statement does not cover \texttt{ABBA2Midpoint}: arithmetic averaging of
the final copies is not a symplectic projection in general.  Nor does it
automatically cover an energy diagnostic obtained by advancing or
reconstructing $k$ after a physical-only projected map.  For a theorem about
$\K$, the accepted map must be symplectic in the complete canonical extended
space, not merely area preserving in $(x,y)$.

\section{Finite regularity of the runtime potential}

The analytic hypothesis is substantive for the current implementation.
A tensor-product spline of degree $q$ is generically only $C^{q-1}$ across
simple internal knots.  It is polynomial, hence analytic, inside one open
spline cell, but a trajectory crossing multiple knots does not remain in one
common analytic expression.  The HDF5 evaluator periodically wraps coordinates
and extends the sampled values across the cell seam.  It therefore introduces
no separate clipping discontinuity; crossings of the periodic seam have the
same finite-smoothness limitation as crossings of ordinary spline knots.

Consequently, Theorem~\ref{thm:no-secular-drift} is a rigorous theorem for an
analytic idealization, but its exponential time scale is not a global theorem
for the periodic-spline runtime representation.  Finite-smoothness
backward-error analysis can still construct a modified Hamiltonian through a
finite order.  If a truncation at order $N>p$ gives a one-step modified-energy
defect $O(h^{N+1})$, the same telescoping argument yields schematically
\begin{equation}
 \abs{\K(Z_n)-\K(Z_0)}
 \leq Ch^p+C t_nh^N.
 \label{eq:finite-smoothness-energy-bound}
\end{equation}
This keeps an $O(h^p)$ band only on an algebraic time scale such as
\begin{equation}
 t_n\leq c h^{-(N-p)},
 \label{eq:algebraic-time-scale}
\end{equation}
and only when the available differentiability and uniform derivative bounds
justify the chosen $N$.

There are two clean ways to strengthen the analytic claim:
\begin{enumerate}
 \item use a globally analytic spatial representation, such as a finite
       Fourier representation on a periodic domain; or
 \item prove a finite-regularity theorem with an explicit $N$ supported by the
       chosen spline degree, including ordinary knots and the periodic seam.
\end{enumerate}

\section{Finite nonlinear solves and floating-point arithmetic}

Let $\PhiNum$ denote the exact-root symplectic map and let the implemented map
be
\begin{equation}
 \widehat\Phi_h(Z)=\PhiNum(Z)+d_h(Z),
 \qquad
 \norm{d_h(Z)}\leq\eta_h.
 \label{eq:perturbed-map}
\end{equation}
Data-dependent Newton or Broyden stopping generally makes
$\widehat\Phi_h$ neither exactly symplectic nor exactly symmetric.  If the
truncated modified Hamiltonian is Lipschitz on the numerical domain, the
additional one-step change caused by $d_h$ is at most $C\eta_h$.  Repeating the
proof of Theorem~\ref{thm:no-secular-drift} gives the worst-case estimate
\begin{equation}
 \abs{\K(\widehat Z_n)-\K(\widehat Z_0)}
 \leq Ch^p+C t_n\exp(-c/h)
       +C\frac{t_n}{h}\eta_h.
 \label{eq:finite-solve-energy-bound}
\end{equation}
The final term is potentially secular.  If the target integration time is
$T$, retaining an $O(h^p)$ energy band in this worst-case estimate requires
\begin{equation}
 \eta_h\ll\frac{h^{p+1}}{T}.
 \label{eq:map-error-requirement}
\end{equation}

The nonlinear residual tolerance is not itself $\eta_h$.  If the root
Jacobian is uniformly invertible and its inverse is bounded by $M$, a local
perturbation argument gives a relation of the form
\begin{equation}
 \eta_h\leq C M\tau_h,
 \label{eq:residual-to-map-error}
\end{equation}
where $\tau_h$ bounds the accepted residual and the constant also depends on
the reconstruction map.  This conditioning factor must be included in any
rigorous implementation-level estimate.

Round-off adds another perturbation.  A random-walk growth proportional to
$\sqrt{n}$ is often observed, but it is a statistical model, not a deterministic
guarantee.  An adversarial deterministic bound must allow growth proportional
to $n$.

\section{What is proved, observed, and not claimed}

\begin{itemize}
 \item \textbf{Proved for the ideal analytic map:} Gauss--Legendre and the
       fully extended implicit ABBA variants have an $O(h^p)$ generalized-energy
       band over exponentially long, but finite, times.
 \item \textbf{Proved algebraically for the exact physical flow:}
       $H(t,z)$ is not generally conserved, whereas $\K=H+k$ is conserved.
 \item \textbf{Not proved by a bounded-looking graph:} sign changes, a small
       final-to-peak ratio, and a small fitted slope support the interpretation
       but do not establish the long-time theorem.
 \item \textbf{Not covered by the exponential theorem as currently represented:}
       spline-knot and periodic-seam crossings, finite nonlinear
       stopping, and floating-point arithmetic.
 \item \textbf{Not claimed:} exact conservation of $\K$, a bound for all
       $n$, or a bounded physical difference $H(t_n,z_n)-H(0,z_0)$.
\end{itemize}

If exact step-by-step generalized-energy conservation is the primary goal, an
energy-preserving discrete-gradient method gives a different algebraic result,
\begin{equation}
 \K(Z_{n+1})=\K(Z_n),
 \label{eq:discrete-gradient-contrast}
\end{equation}
at the exact nonlinear root.  Such a method is not generally symplectic.  This
is stronger in energy and weaker in symplectic structure than the near-
conservation result proved here.

\section{Conclusion}

The absence of secular generalized-energy drift for symplectic GC2D
integrators is a backward-error statement.  Symplecticity creates a modified
Hamiltonian, analyticity makes the defect in its numerical conservation
exponentially small, and comparison of the modified and original Hamiltonians
confines $\K$ to an $O(h^p)$ band for exponentially long times.  Symmetry
improves the parity of the modified expansion but is not, by itself, the
source of the energy bound.

For the project, this theorem applies directly to the exact-root idealizations
of Gauss--Legendre and the fully extended implicit ABBA methods.  A literal
implementation-level theorem additionally requires a regularity statement for
the periodic spline, conditioning bounds for the implicit equations, and
explicit control of nonlinear and round-off errors.

\section*{References}

\begin{itemize}
 \item G. Benettin and A. Giorgilli,
       ``On the Hamiltonian interpolation of near-to-the-identity symplectic
       mappings with application to symplectic integration algorithms,''
       \emph{Journal of Statistical Physics}, 74 (1994), 1117--1143.
 \item E. Hairer and C. Lubich,
       ``The life-span of backward error analysis for numerical integrators,''
       \emph{Numerische Mathematik}, 76 (1997), 441--462.
 \item E. Hairer, C. Lubich, and G. Wanner,
       \emph{Geometric Numerical Integration: Structure-Preserving Algorithms
       for Ordinary Differential Equations}, second edition, Springer, 2006,
       Chapter~IX.
 \item J. M. Sanz-Serna,
       ``Runge--Kutta schemes for Hamiltonian systems,''
       \emph{BIT Numerical Mathematics}, 28 (1988), 877--883.
\end{itemize}

\end{document}
